\section{لایه سخت‌افزاری (بخش پشتیبان امنیت)}
لایه سخت‌افزاری به‌عنوان بخش پشتیبان امنیت، نگهداری اسرار رمزنگاری و امضای پیام‌های خروجی را در دستگاه‌هایی مقاوم در برابر دست‌کاری\LTRfootnote{\lr{Tamper-Proof Device, TPD}} انجام می‌دهد و دست‌کاری سخت‌افزاری یکی از حملاتی است که در مرورهای امنیتی ارتباطات خودرویی، ساختارهای سخت‌افزاری مورد اعتماد برای مقابله با آن معرفی شده‌اند \cite{SecuringVehicularAd, farsimadanReviewSecurityChallenges2025}. مرور سیستماتیک مطالعات حفظ حریم خصوصی در شبکه‌های اقتضایی خودرویی نیز تقویت واحد سوارشونده با سخت‌افزار مورد اعتماد را در میان راه‌حل‌های بررسی‌شده فهرست کرده است \cite{rahmawatiagustinaSystematicLiteratureReview2025}.

\subsection{ماژول پلتفرم قابل اعتماد}
در یکی از معماری‌های بنیادین امنیت شبکهٔ اقتضایی خودرویی، دستگاه مقاوم در برابر دست‌کاری اسرار رمزنگاری را نگه می‌دارد و پیام‌های خروجی را امضا می‌کند؛ این دستگاه به باتری و ساعت مستقلی مجهز است که ساعت آن هنگام عبور از ایستگاه‌های پایهٔ مورد اعتماد هم‌گام‌سازی می‌شود و حسگرهای دست‌کاری آن در صورت دست‌کاری تمام کلیدها را پاک می‌کنند \cite{SecuringVehicularAd}. نویسندگان همان مقاله ضعف این دستگاه را در حساسیت بیش از حد به شرایط جاده (ضربه و دمای شدید) و گرانی آن دانسته‌اند و نمونهٔ \lr{IBM 4758} را با هزینهٔ چند هزار دلار آورده‌اند؛ در مقابل، ماژول سکوی مورداعتماد\LTRfootnote{\lr{Trusted Platform Module, TPM}} را با هزینه‌ای به اندازهٔ چند ده دلار معرفی کرده‌اند که در برابر حملات نرم‌افزاری مقاوم است، نه در برابر دست‌کاری سخت‌افزاری، و در سال ۲۰۰۷ پیش‌بینی کرده بودند سخت‌افزار نهایی احتمالاً سازشی میان این دو خواهد بود \cite{SecuringVehicularAd}. مرور چالش‌های امنیتی ارتباطات \lr{V2X} نیز دست‌کاری سخت‌افزاری هنگام سرویس سالانهٔ خودرو را در دستهٔ حملات بر زیرساخت آورده و ماژول سکوی مورداعتماد را دفاعِ آن معرفی کرده است \cite{farsimadanReviewSecurityChallenges2025}. مقالهٔ بنیادین، مقاومت ماژول سکوی مورداعتماد را در برابر دست‌کاری سخت‌افزاری نمی‌داند و آن را تنها در برابر حملات نرم‌افزاری مؤثر می‌شمارد؛ بنابراین میان این دو منبع دربارهٔ مقابله با این حمله تنشی وجود دارد \cite{SecuringVehicularAd}.

\subsection{توابع فیزیکی غیرقابل شبیه‌سازی}
فراتر از دستگاه مقاوم در برابر دست‌کاری و ماژول سکوی مورداعتماد، مرور سیستماتیکِ ۱۱۳ مطالعه (۲۰۱۹ تا ۲۰۲۴) تابع غیرقابل همسان‌سازی فیزیکی\LTRfootnote{\lr{Physically Unclonable Function, PUF}} را نیز در کنار دستگاه مقاوم در برابر دست‌کاری و بیومتریک، از جمله وسایل تقویت واحد سوارشونده در مطالعات بررسی‌شده دانسته است \cite{rahmawatiagustinaSystematicLiteratureReview2025}. جزئیات فنی عملکرد این توابع در سه منبع این فصل شرح داده نشده و بنابراین بیش از همان ذکر صریح در اینجا امکان‌پذیر نیست.

\subsection{تصدیق چندعاملی}
ترکیب این سخت‌افزارها با عوامل دیگر احراز هویت، یعنی تصدیق چندعاملی\LTRfootnote{\lr{Multi-Factor Authentication, MFA}}، در منابع این فصل پوشش داده نشده است؛ تنها بیومتریک در فهرست وسایل تقویت واحد سوارشوندهٔ مرور سیستماتیک آمده و طرحی چندعاملی در هیچ‌یک از این منابع معرفی نشده است \cite{rahmawatiagustinaSystematicLiteratureReview2025}.

\subsection{مقابله با حملات کانال جانبی}
همچون تصدیق چندعاملی، حملات کانال جانبی\LTRfootnote{\lr{Side-Channel Attacks}} و تکنیک‌های مقابله با آن نیز در هیچ‌یک از سه منبع مرورشدهٔ این فصل وارد نشده و این زیربخش تنها به‌عنوان پیوندی به موضوعی مجاور حفظ شده است.

\subsection{عنصر امن سخت‌افزاری و محیط اجرای مطمئن}
عنصر امن سخت‌افزاری\LTRfootnote{\lr{Embedded Secure Element, eSE}} و محیط اجرای مطمئن\LTRfootnote{\lr{Trusted Execution Environment, TEE}} نیز از دیگر ساختارهای سخت‌افزاری امنیتی‌اند که در هیچ‌یک از منابع این فصل پوشش داده نشده‌اند و بنابراین هیچ ادعای فنی دربارهٔ آن‌ها در اینجا ثبت نمی‌شود.

\subsection{ملاحظات هزینه و مدیریت}
ارزیابی این سخت‌افزارها را نمی‌توان از هزینه و مدیریت چرخهٔ عمر آن‌ها جدا کرد. نویسندگان مقالهٔ بنیادین گرانی دستگاه مقاوم در برابر دست‌کاری (نمونهٔ \lr{IBM 4758} با چند هزار دلار در برابر چند ده دلار ماژول سکوی مورداعتماد) و حساسیت آن به شرایط جاده را از نقاط ضعف آن دانسته و سازش میان این دو را در سال ۲۰۰۷ پیش‌بینی کرده بودند \cite{SecuringVehicularAd}. مرور چالش‌های امنیتی نیز هزینهٔ فناوری‌های \lr{V2X} را از عوامل افزایش قیمت خودرو و پراکندگی زیرساخت کنارجاده‌ای به دلیل هزینهٔ نصب و بهره‌بردار دانسته است؛ به‌روزرسانی‌های از راه دور را نیز به‌تدریج استاندارد می‌داند، اما خودروهایی که اتصالشان قطع شده را در برابر به‌روزرسانی‌های ناقص یا از دست‌رفته آسیب‌پذیر می‌شمارد \cite{farsimadanReviewSecurityChallenges2025}. چرخهٔ عمر کلیدها شامل بارگذاری کلیدها هنگام ثبت خودرو، تجدید دوره‌ای کلیدهای ناشناس (مثلاً در معاینهٔ سالانه) و پروتکل‌های ابطال \lr{RTPD}، \lr{RCCRL} و \lr{DRP} (پیام رمزشدهٔ مرجع صدور گواهی به دستگاه مقاوم در برابر دست‌کاری برای پاک‌سازی کلیدها، فهرست ابطال فشرده و انباشت اتهام توسط همسایه‌ها) است و ابطال را احتمالاً دشوارترین جنبهٔ امنیت این شبکه‌ها می‌داند \cite{SecuringVehicularAd}. از این‌رو، تعادل میان امنیت و هزینه باید در طراحی سامانه‌های خودرویی لحاظ شود.
